\chapter{2019年全国I卷（文）}

\section{单选题}

\begin{problem}
设 \( z=\dfrac{3-\i}{1+2\i} \)，则 \( |z|=\)（\quad）

\choices
    {\( 2 \)}
    {\( \sqrt3 \)}
    {\( \sqrt2 \)}
    {\( 1 \)}
\end{problem}

\begin{answer}
C
\end{answer}

\begin{problem}
已知集合 \( U=\{1,2,3,4,5,6,7\} \)，\( A=\{2,3,4,5\} \)，\( B=\{2,3,6,7\} \)，则 \( B\cap \complement_U A=\)（\quad）

\choices
    {\( \{1,6\} \)}
    {\( \{1,7\} \)}
    {\( \{6,7\} \)}
    {\( \{1,6,7\} \)}
\end{problem}

\begin{answer}
C
\end{answer}

\begin{problem}
已知 \( a=\log_2 0.2 \)，\( b=2^{0.2} \)，\( c=0.2^{0.3} \)，则（\quad）

\choices
    {\( a<b<c \)}
    {\( a<c<b \)}
    {\( c<a<b \)}
    {\( b<c<a \)}
\end{problem}

\begin{answer}
B
\end{answer}

\begin{problem}
古希腊时期，人们认为最美人体的头顶至肚脐的长度与肚脐至足底的长度之比是 \( \dfrac{\sqrt5-1}{2} \)（\( \dfrac{\sqrt5-1}{2}\approx 0.618 \)，称为黄金分割比例），著名的“断臂维纳斯”便是如此. 此外，最美人体的头顶至咽喉的长度与咽喉至肚脐的长度之比也是 \( \dfrac{\sqrt5-1}{2} \). 若某人满足上述两个黄金分割比例，且腿长为 \( 105\text{ cm} \)，头顶至脖子下端的长度为 \( 26\text{ cm} \)，则其身高可能是（\quad）

\bitmapfigure[max width=6.0cm]{img/2019/national_paper_1_liberal/q4_fig1.png}

\choices
    {\( 165\text{ cm} \)}
    {\( 175\text{ cm} \)}
    {\( 185\text{ cm} \)}
    {\( 190\text{ cm} \)}
\end{problem}

\begin{answer}
B
\end{answer}

\begin{problem}
函数 \( f(x)=\dfrac{\sin x+x}{\cos x+x^2} \) 在 \( [-\pi,\pi] \) 的图象大致为（\quad）

\choices
    {\(\choicebitmap[width=0.15\paperwidth]{img/2019/national_paper_1_liberal/q5_choice_a.png}\)}
    {\(\choicebitmap[width=0.15\paperwidth]{img/2019/national_paper_1_liberal/q5_choice_b.png}\)}
    {\(\choicebitmap[width=0.15\paperwidth]{img/2019/national_paper_1_liberal/q5_choice_c.png}\)}
    {\(\choicebitmap[width=0.15\paperwidth]{img/2019/national_paper_1_liberal/q5_choice_d.png}\)}
\end{problem}

\begin{answer}
D
\end{answer}

\begin{problem}
某学校为了了解 \( 1000 \) 名新生的身体素质，将这些学生编号 \(1\)，\(2\)，\(\dots\)，\(1000\)，从这些新生中用系统抽样方法等距抽取 \( 100 \) 名学生进行体质测验. 若 \( 46 \) 号学生被抽到，则下面 \( 4 \) 名学生中被抽到的是（\quad）

\choices
    {\( 8 \) 号学生}
    {\( 200 \) 号学生}
    {\( 616 \) 号学生}
    {\( 815 \) 号学生}
\end{problem}

\begin{answer}
C
\end{answer}

\begin{problem}
\( \tan 255^\circ=\)（\quad）

\choices
    {\( -2-\sqrt3 \)}
    {\( -2+\sqrt3 \)}
    {\( 2-\sqrt3 \)}
    {\( 2+\sqrt3 \)}
\end{problem}

\begin{answer}
D
\end{answer}

\begin{problem}
已知非零向量 \(\bs{a}\)，\(\bs{b}\) 满足 \( |\bs{a}|=2|\bs{b}| \)，且 \( (\bs{a}-\bs{b})\perp \bs{b} \)，则 \( \bs{a} \) 与 \( \bs{b} \) 的夹角为（\quad）

\choices
    {\( \dfrac{\pi}{6} \)}
    {\( \dfrac{\pi}{3} \)}
    {\( \dfrac{2\pi}{3} \)}
    {\( \dfrac{5\pi}{6} \)}
\end{problem}

\begin{answer}
B
\end{answer}

\begin{problem}
如图是求 \( \dfrac{1}{2+\dfrac{1}{2+\dfrac{1}{2}}} \) 的程序框图，图中空白框中应填入（\quad）

\texfigure[max width=0.32\linewidth]{img_repaint/2019/national_paper_1_liberal/q9_fig1.tex}

\choices
    {\( A=\dfrac{1}{2+A} \)}
    {\( A=2+\dfrac{1}{A} \)}
    {\( A=\dfrac{1}{1+2A} \)}
    {\( A=1+\dfrac{1}{2A} \)}
\end{problem}

\begin{answer}
A
\end{answer}

\begin{problem}
双曲线 \( C:\dfrac{x^2}{a^2}-\dfrac{y^2}{b^2}=1 \)（\( a>0 \)，\( b>0 \)）的一条渐近线的倾斜角为 \( 130^\circ \)，则 \(C\) 的离心率为（\quad）

\choices
    {\( 2\sin 40^\circ \)}
    {\( 2\cos 40^\circ \)}
    {\( \dfrac{1}{\sin 50^\circ} \)}
    {\( \dfrac{1}{\cos 50^\circ} \)}
\end{problem}

\begin{answer}
D
\end{answer}

\begin{problem}
\( \triangle ABC \) 的内角 \(A\)，\(B\)，\(C\) 的对边分别为 \(a\)，\(b\)，\(c\). 已知 \( a\sin A-b\sin B=4c\sin C \)，\( \cos A=-\dfrac{1}{4} \)，则 \( \dfrac{b}{c}=\)（\quad）

\choices
    {\( 6 \)}
    {\( 5 \)}
    {\( 4 \)}
    {\( 3 \)}
\end{problem}

\begin{answer}
A
\end{answer}

\begin{problem}
已知椭圆 \(C\) 的焦点为 \( F_1(-1,0) \)，\( F_2(1,0) \)，过 \(F_2\) 的直线与 \(C\) 交于 \(A\)，\(B\) 两点. 若 \( |AF_2|=2|F_2B| \)，\( |AB|=|BF_1| \)，则 \(C\) 的方程为（\quad）

\choices
    {\( \dfrac{x^2}{2}+y^2=1 \)}
    {\( \dfrac{x^2}{3}+\dfrac{y^2}{2}=1 \)}
    {\( \dfrac{x^2}{4}+\dfrac{y^2}{3}=1 \)}
    {\( \dfrac{x^2}{5}+\dfrac{y^2}{4}=1 \)}
\end{problem}

\begin{answer}
B
\end{answer}

\section{填空题}

\begin{problem}
曲线 \( y=3(x^2+x)\e^x \) 在点 \( (0,0) \) 处的切线方程为\fillinblank{}.
\end{problem}

\begin{answer}
\(y=3x\)
\end{answer}

\begin{problem}
记 \( S_n \) 为等比数列 \( \{a_n\} \) 的前 \( n \) 项和. 若 \( a_1=1 \)，\( S_3=\dfrac{3}{4} \)，则 \( S_4=\)\fillinblank{}.
\end{problem}

\begin{answer}
\(\dfrac{5}{8}\)
\end{answer}

\begin{problem}
函数 \( f(x)=\sin \left( 2x+\dfrac{3\pi}{2} \right)-3\cos x \) 的最小值为\fillinblank{}.
\end{problem}

\begin{answer}
\(-4\)
\end{answer}

\begin{problem}
已知 \( \angle ACB=90^\circ \)，\(P\) 为平面 \(ABC\) 外一点，\( PC=2 \)，点 \(P\) 到 \( \angle ACB \) 两边 \(AC\)，\(BC\) 的距离均为 \( \sqrt3 \)，那么 \(P\) 到平面 \(ABC\) 的距离为\fillinblank{}.
\end{problem}

\begin{answer}
\(\sqrt2\)
\end{answer}

\section{解答题}

\begin{problem}
某商场为提高服务质量，随机调查了 \( 50 \) 名男顾客和 \( 50 \) 名女顾客，每位顾客对该商场的服务给出满意或不满意的评价，得到下面列联表：

\begin{center}
\begin{tabular}{c|cc}
 & 满意 & 不满意 \\ \hline
男顾客 & 40 & 10 \\
女顾客 & 30 & 20
\end{tabular}
\end{center}

\begin{enumerate}
\item 分别估计男、女顾客对该商场服务满意的概率；
\item 能否有 \( 95\% \) 的把握认为男、女顾客对该商场服务的评价有差异？
\end{enumerate}

附：\( K^2=\dfrac{n(ad-bc)^2}{(a+b)(c+d)(a+c)(b+d)} \).

\begin{center}
\begin{tabular}{c|ccc}
\( P(K^2\ge k) \) & \( 0.050 \) & \( 0.010 \) & \( 0.001 \) \\ \hline
\( k \) & \( 3.841 \) & \( 6.635 \) & \( 10.828 \)
\end{tabular}
\end{center}
\end{problem}

\begin{solution}
\begin{enumerate}
\item 男顾客中对服务满意的有 \(40\) 人，共调查 \(50\) 人，所以男顾客对服务满意的概率估计为
\[
\frac{40}{50}=0.8
\]
女顾客中对服务满意的有 \(30\) 人，共调查 \(50\) 人，所以女顾客对服务满意的概率估计为
\[
\frac{30}{50}=0.6
\]

\item 在题给公式中取 \(a=40\)，\(b=10\)，\(c=30\)，\(d=20\)，\(n=100\)，得
\[
K^2=\frac{100(40\times20-10\times30)^2}{50\times50\times70\times30}=\frac{100}{21}\approx4.762
\]
由于 \(4.762>3.841\)，所以有 \(95\%\) 的把握认为男、女顾客对该商场服务的评价有差异.
\end{enumerate}
\end{solution}

\begin{problem}
记 \( S_n \) 为等差数列 \( \{a_n\} \) 的前 \( n \) 项和. 已知 \( S_9=-a_5 \).

\begin{enumerate}
\item 若 \( a_3=4 \)，求 \( \{a_n\} \) 的通项公式；
\item 若 \( a_1>0 \)，求使得 \( S_n\ge a_n \) 的 \( n \) 的取值范围.
\end{enumerate}
\end{problem}

\begin{solution}
\begin{enumerate}
\item 设等差数列 \(\{a_n\}\) 的首项为 \(a_1\)，公差为 \(d\). 由等差数列的前 \(n\) 项和公式，
\[
S_9=\frac{9}{2}(2a_1+8d)=9a_1+36d
\]
且 \(a_5=a_1+4d\). 代入已知条件 \(S_9=-a_5\)，得
\[
9a_1+36d=-(a_1+4d)
\]
即 \(a_1+4d=0\). 又因为 \(a_3=a_1+2d=4\)，所以
\(a_1=-4d\)，\(-4d+2d=4\)
解得 \(d=-2\)，\(a_1=8\). 因而
\[
a_n=a_1+(n-1)d=8-2(n-1)=10-2n
\]

\item 由 \(a_1+4d=0\) 和 \(a_1>0\)，得 \(d<0\). 于是
\(S_n=\frac{n}{2}\bigl(2a_1+(n-1)d\bigr)=\frac{d}{2}n(n-9)\)，\(a_n=a_1+(n-1)d=(n-5)d\).
因此
\[
S_n\ge a_n
\iff \frac{d}{2}n(n-9)\ge d(n-5)
\]
由于 \(d<0\)，不等式同除以 \(d/2\) 时不等号改变方向，故
\[
n(n-9)\le2(n-5)
\iff n^2-11n+10\le0
\iff (n-1)(n-10)\le0
\]
又 \(n\in\mathbb N^*\)，所以 \(n\) 的取值范围为 \(1\le n\le10\).
\end{enumerate}
\end{solution}

\begin{problem}
如图，直四棱柱 \( ABCD-A_1B_1C_1D_1 \) 的底面是菱形，\( AA_1=4 \)，\( AB=2 \)，\( \angle BAD=60^\circ \)，\(E\)，\(M\)，\(N\) 分别是 \(BC\)，\(BB_1\)，\(A_1D\) 的中点.

\begin{enumerate}
\item 证明：\( MN\parallel \) 平面 \(C_1DE\)；
\item 求点 \(C\) 到平面 \(C_1DE\) 的距离.
\end{enumerate}

\bitmapfigure[max width=0.26\linewidth]{img/2019/national_paper_1_liberal/q19_fig1.png}
\end{problem}

\begin{solution}
\begin{enumerate}
\item 建立空间直角坐标系，使
\(A=(0,0,0)\)，\(B=(2,0,0)\)，\(D=(1,\sqrt3,0)\)，\(C=(3,\sqrt3,0)\).
因为四棱柱为直四棱柱且 \(AA_1=4\)，所以
\(C_1=(3,\sqrt3,4)\)，\(E=\left(\frac52,\frac{\sqrt3}{2},0\right)\)，\(M=(2,0,2)\)，\(N=\left(\frac12,\frac{\sqrt3}{2},2\right)\).
于是
\(\overrightarrow{DE}=\left(\frac32,-\frac{\sqrt3}{2},0\right)\)，\(\overrightarrow{MN}=\left(-\frac32,\frac{\sqrt3}{2},0\right)=-\overrightarrow{DE}\).
平面 \(C_1DE\) 的一个法向量为 \((2,2\sqrt3,-1)\)，故其方程为
\[
2x+2\sqrt3y-z-8=0
\]
将 \(M\)，\(N\) 的坐标代入左边，均得 \(-6\)，所以直线 \(MN\) 位于平面 \(2x+2\sqrt3y-z-2=0\) 内，该平面与平面 \(C_1DE\) 平行. 因而 \(MN\parallel\) 平面 \(C_1DE\).

\item 点 \(C=(3,\sqrt3,0)\)，所以点 \(C\) 到平面 \(C_1DE\) 的距离为
\[
\frac{|2\times3+2\sqrt3\times\sqrt3-0-8|}{\sqrt{2^2+(2\sqrt3)^2+(-1)^2}}
=\frac4{\sqrt{17}}
\]
\end{enumerate}
\end{solution}

\begin{problem}
已知函数 \( f(x)=2\sin x-x\cos x-x \)，\( f'(x) \) 为 \( f(x) \) 的导数.

\begin{enumerate}
\item 证明：\( f'(x) \) 在区间 \( (0,\pi) \) 存在唯一零点；
\item 若 \( x\in [0,\pi] \) 时，\( f(x)\ge ax \)，求 \( a \) 的取值范围.
\end{enumerate}
\end{problem}

\begin{solution}
\begin{enumerate}
\item 求导得
\[
f'(x)=2\cos x-\bigl(\cos x-x\sin x\bigr)-1=x\sin x+\cos x-1
\]
令 \(g(x)=f'(x)\)，则
\[
g'(x)=x\cos x
\]
在 \(\left(0,\frac\pi2\right)\) 上，\(g'(x)>0\)，且 \(g(0)=0\)，所以 \(g(x)>0\). 在 \(\left(\frac\pi2,\pi\right)\) 上，\(g'(x)<0\)，又
\(g\left(\frac\pi2\right)=\frac\pi2-1>0\)，\(g(\pi)=-2<0\).
由连续函数零点存在定理，\(g(x)\) 在 \(\left(\frac\pi2,\pi\right)\) 上至少有一个零点；由其在该区间上严格递减，零点至多有一个. 因此 \(f'(x)\) 在 \((0,\pi)\) 上存在唯一零点.

\item 对 \(x\in[0,\pi]\)，有
\[
f(x)=2\sin x-x(1+\cos x)
\]
当 \(0\le x<\pi\) 时，令 \(u=\dfrac{x}{2}\)，则
\[
f(x)=2\cos u\bigl(2\sin u-x\cos u\bigr)=4\cos^2u\bigl(\tan u-u\bigr)
\]
又因为 \(0\le u<\dfrac\pi2\)，且
\((\tan u-u)'=\tan^2u\ge0\)，\(\tan0-0=0\)
所以 \(\tan u-u\ge0\)，从而 \(f(x)\ge0\). 当 \(x=\pi\) 时，\(f(\pi)=0\)，故在整个区间 \([0,\pi]\) 上均有 \(f(x)\ge0\).

若 \(a\le0\)，则因 \(x\ge0\)，有 \(ax\le0\)，从而 \(f(x)\ge0\ge ax\). 反之，令 \(x=\pi\)，由 \(f(\pi)=0\ge a\pi\) 可得 \(a\le0\). 因此 \(a\) 的取值范围为 \((-\infty,0]\).
\end{enumerate}
\end{solution}

\begin{problem}
已知点 \(A\)，\(B\) 关于坐标原点 \(O\) 对称，\( |AB|=4 \)，\( \odot M \) 过点 \(A\)，\(B\) 且与直线 \( x+2=0 \) 相切.

\begin{enumerate}
\item 若 \(A\) 在直线 \( x+y=0 \) 上，求 \( \odot M \) 的半径；
\item 是否存在定点 \(P\)，使得当 \(A\) 运动时，\( |MA|-|MP| \) 为定值？并说明理由.
\end{enumerate}
\end{problem}

\begin{solution}
设 \(A=(x,y)\)，则 \(B=(-x,-y)\)，且 \(x^2+y^2=4\). 设圆心 \(M=(u,v)\)，圆的半径为 \(r\). 因为圆经过 \(A\)，\(B\)，所以
\[
(u-x)^2+(v-y)^2=(u+x)^2+(v+y)^2
\]
即 \(ux+vy=0\). 因而
\[
r^2=MA^2=u^2+v^2+x^2+y^2=u^2+v^2+4
\]
又因为圆与直线 \(x+2=0\) 相切，圆心到该直线的距离等于半径，所以
\[
(u+2)^2=r^2=u^2+v^2+4
\]
从而
\[
v^2=4u
\]
特别地，\(u\ge0\)，且 \(r=u+2\).

\begin{enumerate}
\item 由 \(A\) 在直线 \(x+y=0\) 上，可设 \(A=(s,-s)\). 由 \(x^2+y^2=4\) 得 \(s=\pm\sqrt2\). 代入 \(ux+vy=0\)，得 \(u=v\). 再结合 \(v^2=4u\)，有
\[
u^2=4u
\]
所以 \(u=0\) 或 \(u=4\). 因而圆的半径为
\[
r=u+2=2\quad\text{或}\quad6
\]

\item 取定点 \(P=(1,0)\). 由 \(v^2=4u\)，有
\[
MP^2=(u-1)^2+v^2=(u-1)^2+4u=(u+1)^2
\]
由于 \(u\ge0\)，所以 \(MP=u+1\). 结合 \(MA=r=u+2\)，得到
\[
|MA|-|MP|=(u+2)-(u+1)=1
\]
这个差与点 \(A\) 的位置无关，因此存在定点 \(P=(1,0)\)，且定值为 \(1\).
\end{enumerate}
\end{solution}

\section{选考题：解答题}

\begin{problem}
在直角坐标系 \( xOy \) 中，曲线 \(C\) 的参数方程为 \( x=\dfrac{1-t^2}{1+t^2} \)，\( y=\dfrac{4t}{1+t^2} \)（\( t \) 为参数），以坐标原点 \(O\) 为极点，\( x \) 轴的正半轴为极轴建立极坐标系，直线 \( l \) 的极坐标方程为 \( 2\rho \cos \theta+\sqrt3\,\rho \sin \theta+11=0 \).

\begin{enumerate}
\item 求 \(C\) 和 \( l \) 的直角坐标方程；
\item 求 \(C\) 上的点到 \( l \) 距离的最小值.
\end{enumerate}
\end{problem}

\begin{solution}
\begin{enumerate}
\item 由参数方程得
\[
x^2+\frac{y^2}{4}=\frac{(1-t^2)^2}{(1+t^2)^2}+\frac{4t^2}{(1+t^2)^2}=1
\]
所以曲线 \(C\) 的直角坐标方程为
\[
x^2+\frac{y^2}{4}=1
\]
由极坐标与直角坐标的关系 \(x=\rho\cos\theta\)，\(y=\rho\sin\theta\)，直线 \(l\) 的方程化为
\[
2x+\sqrt3y+11=0
\]

\item 在椭圆 \(C\) 上设 \(x=\cos\varphi\)，\(y=2\sin\varphi\). 则
\[
2x+\sqrt3y=2\cos\varphi+2\sqrt3\sin\varphi
\]
由柯西不等式，
\[
2\cos\varphi+2\sqrt3\sin\varphi\ge-\sqrt{2^2+(2\sqrt3)^2}=-4
\]
因此对 \(C\) 上任意一点，都有 \(2x+\sqrt3y+11\ge7>0\). 该点到直线 \(l\) 的距离为
\[
d=\frac{2x+\sqrt3y+11}{\sqrt{2^2+(\sqrt3)^2}}\ge\frac7{\sqrt7}=\sqrt7
\]
当 \((\cos\varphi,\sin\varphi)=\left(-\frac12,-\frac{\sqrt3}{2}\right)\) 时取等号，所以 \(C\) 上的点到 \(l\) 距离的最小值为 \(\sqrt7\).
\end{enumerate}
\end{solution}

\begin{problem}
已知 \(a\)，\(b\)，\(c\) 为正数，且满足 \( abc=1 \). 证明：

\begin{enumerate}
\item \( \dfrac{1}{a}+\dfrac{1}{b}+\dfrac{1}{c}\le a^2+b^2+c^2 \)；
\item \( (a+b)^3+(b+c)^3+(c+a)^3\ge 24 \).
\end{enumerate}
\end{problem}

\begin{solution}
\begin{enumerate}
\item 因为 \(abc=1\)，所以
\[
\frac1a+\frac1b+\frac1c=bc+ca+ab
\]
又
\[
a^2+b^2+c^2-(ab+bc+ca)
=\frac12\bigl((a-b)^2+(b-c)^2+(c-a)^2\bigr)\ge0
\]
因此
\[
\frac1a+\frac1b+\frac1c\le a^2+b^2+c^2
\]

\item 由基本不等式，
\[
(a+b)^3\ge\bigl(2\sqrt{ab}\bigr)^3=8(ab)^{3/2}
\]
同理可得
\((b+c)^3\ge8(bc)^{3/2}\)，\((c+a)^3\ge8(ca)^{3/2}\).
相加得
\[
(a+b)^3+(b+c)^3+(c+a)^3
\ge8\bigl((ab)^{3/2}+(bc)^{3/2}+(ca)^{3/2}\bigr)
\]
再次使用基本不等式，得
\[
(ab)^{3/2}+(bc)^{3/2}+(ca)^{3/2}
\ge3\sqrt[3]{(ab)^{3/2}(bc)^{3/2}(ca)^{3/2}}
=3\sqrt[3]{(abc)^3}=3
\]
所以
\[
(a+b)^3+(b+c)^3+(c+a)^3\ge8\times3=24
\]
\end{enumerate}
\end{solution}
